% Kandydat restrukturyzacji; nie jest aktywnym rozdziałem Atlasu.
\section{Dystrybucje o zwartym nośniku}
\providecommand{\atlasref}[2]{#2}

Pracujemy nad $\K\in\{\R,\C\}$. Przez $\mathcal E'(\Omega)$ rozumiemy przestrzeń sprzężoną do $C^\infty(\Omega)$; jej identyfikację z dystrybucjami o zwartym nośniku udowodnimy poniżej.

\subsection{Dystrybucje o zwartym nośniku i funkcje gładkie}

\begin{lemat}[Przedłużenie dystrybucji o zwartym nośniku]
\label{lem:compact-action21}
Niech $T\in\mathcal D'(\R^d;\K)$ i
$K=\supp T\Subset\R^d$. Wówczas dla każdej funkcji
$f\in C^\infty(\R^d;\K)$ istnieje dobrze określona wartość
$\langle T,f\rangle:=\langle T,\chi f\rangle$, gdzie
$\chi\in\mathcal D(\R^d)$ jest równa $1$ na otoczeniu $K$.
Wartość ta nie zależy od wyboru $\chi$. Istnieje także
niepusty zbiór zwarty $L$ i liczba $m\in\N_0$ oraz $C>0$ takie, że
\[
 |\langle T,f\rangle|\le
 C\max_{|\alpha|\le m}\sup_{x\in L}|\partial^\alpha f(x)|
 \quad(f\in C^\infty(\R^d)).
\]
W szczególności $T$ wyznacza ciągły funkcjonał na
przestrzeni Frécheta $C^\infty(\R^d)$.
\end{lemat}
\begin{proof}
Funkcja $\chi$ istnieje na mocy Lematu~\ref{lem:cutoffs17}.
Jeżeli $\chi_1,\chi_2$ są równe $1$ w otoczeniach $K$,
to $(\chi_1-\chi_2)f$ zeruje się na otoczeniu $K$;
jej nośnik jest rozłączny z $K$, więc z Lematu~\ref{lem:distribution-local19}
wynika $\langle T,(\chi_1-\chi_2)f\rangle=0$.
Gdy $f\in\mathcal D(\R^d)$, identyczny argument dla
$(1-\chi)f$ pokazuje, że przedłużenie zgadza się z pierwotnym
działaniem $T$. Liniowość wynika z użycia jednej ustalonej $\chi$.
Dla ustalonej $\chi$ wybieramy niepusty zwarty zbiór $L$
zawierający $\supp\chi$.
Ciągłość $T$ na $\mathcal D_L$ daje nierówność
$|\langle T,\psi\rangle|\le C_0p_m(\psi)$ dla funkcji
$\psi\in\mathcal D_L$. Podstawiamy $\psi=\chi f$
i stosujemy wzór Leibniza. Dokładniej, dla $|\alpha|\le m$ mamy
$\partial^\alpha(\chi f)=\sum_{\beta\le\alpha}
\binom\alpha\beta(\partial^\beta\chi)(\partial^{\alpha-\beta}f)$.
Poza $L$ wszystkie składniki znikają. Na $L$ ich moduły
ograniczamy przez iloczyny odpowiednich supremów.
Możemy zatem przyjąć
\[
 C=C_0\max\left\{1,\max_{|\alpha|\le m}
 \sum_{\beta\le\alpha}\binom\alpha\beta
 \sup_{x\in\R^d}|\partial^\beta\chi(x)|\right\}.
\]
Otrzymujemy żądane oszacowanie, a zatem ciągłość na
$C^\infty(\R^d)$ także w przypadku $T=0$.
\end{proof}

\begin{uwaga}[Przestrzeń dualna przestrzeni $C^\infty$]
\label{rem:Eprime21}
Odwrotnie, niech $F$ będzie ciągłym funkcjonałem liniowym na
$C^\infty(\R^d;\K)$. Ciągłość daje oszacowanie
$|F(f)|\le C\max_{|\alpha|\le m}\sup_{x\in L}|\partial^\alpha f(x)|$
dla pewnego zwartego $L\subseteq\R^d$, $m\in\N_0$ i $C>0$:
wystarczy zastąpić skończoną rodzinę seminorm jedną seminormą,
biorąc unię odpowiednich zbiorów zwartych i maksimum rzędów
(oraz w razie potrzeby powiększając $L$ do zbioru niepustego).
Na każdym etapie $\mathcal D_M$ prawa strona jest ograniczona
przez $Cp_m(f)$. Własność uniwersalna granicy induktywnej
(Twierdzenie~\ref{thm:LF-universal18}) daje ciągłość ograniczenia
$T=F|_{\mathcal D}$, a znikanie $F$ na funkcjach testowych
o nośniku rozłącznym z $L$ daje $\supp T\subseteq L$.
Wybierając $\chi\in\mathcal D(\R^d)$ równą $1$ w otoczeniu $L$, mamy
$F(f-\chi f)=0$ z oszacowania, zatem $F(f)=\langle T,\chi f\rangle$.
To pokazuje również jednoznaczność ciągłego przedłużenia.
Otrzymujemy naturalną identyfikację liniową
$\mathcal E'(\R^d)\cong C^\infty(\R^d)'$ przez ograniczenie
i przedłużenie, przy czym poniżej
nie będziemy utożsamiać bez wyjaśnienia różnych możliwych
topologii na tej przestrzeni dualnej.
\end{uwaga}

\subsection{Przestrzeń sprzężona do funkcji gładkich na zbiorze otwartym}
\begin{twierdzenie}[Dystrybucje o zwartym nośniku]
\label{thm:compact-support-E25}
Niech $\Omega\subseteq\R^d$ będzie niepustym zbiorem otwartym.
Obraz kanonicznego odwzorowania
$\mathcal E'(\Omega)\to\mathcal D'(\Omega)$
składa się dokładnie z dystrybucji o zwartym nośniku.
Dokładniej, dla każdego $F\in\mathcal E'(\Omega)$ istnieją
niepusty zwarty $K\Subset\Omega$, $m\in\N_0$ i $C>0$ takie, że
\[
 |F(f)|\le C\max_{|\alpha|\le m}\sup_{x\in K}
 |\partial^\alpha f(x)|\qquad(f\in\mathcal E(\Omega)).
\]
Odwrotnie, jeżeli $T\in\mathcal D'(\Omega)$ ma zwarty nośnik,
to istnieje dokładnie jeden funkcjonał ciągły
$F\in\mathcal E'(\Omega)$ równy $T$ na $\mathcal D(\Omega)$.
\end{twierdzenie}
\begin{proof}
Z kryterium ciągłości funkcjonału na przestrzeni
$\mathcal E(\Omega)$ (Twierdzenie~\ref{thm:dual-cinf17})
otrzymujemy oszacowanie przez skończenie wiele
seminorm $p_{K_j,m_j}$. Biorąc $K=\bigcup_jK_j$
i $m=\max_jm_j$, otrzymujemy oszacowanie z tezy
(dla funkcjonału zerowego można wybrać dowolny
niepusty zwarty $K$). Ograniczenie
$T=F|_{\mathcal D(\Omega)}$ jest dystrybucją, ponieważ
odwzorowanie $\mathcal D(\Omega)\to\mathcal E(\Omega)$
jest ciągłe na mocy twierdzenia~\ref{thm:zanurzenia25}. Gdy
$\varphi\in\mathcal D(\Omega\setminus K)$, jej
przedłużenie przez zero wraz ze wszystkimi pochodnymi
znika w otoczeniu $K$. Zatem $T(\varphi)=0$,
czyli $\supp T\subseteq K$.
Nośnik jest tu nośnikiem dystrybucji $T$,
a nie zbiorem punktów przestrzeni argumentów funkcjonału $F$.

Niech teraz $K=\supp T\Subset\Omega$. Jeśli $K$ jest
pusty, to $T=0$ z lokalności dystrybucji
(Lemat~\ref{lem:distribution-local19}) i wystarcza $F=0$.
W przeciwnym przypadku wybierzmy
$\chi\in\mathcal D(\Omega)$ równą $1$ w otoczeniu $K$
(Lemat~\ref{lem:cutoffs17}) i połóżmy $L=\supp\chi$.
Kładziemy $F(f)=\langle T,\chi f\rangle$ dla
$f\in\mathcal E(\Omega)$. Ponieważ $\supp(\chi f)\subseteq L$,
z lokalnego oszacowania dystrybucji
(Twierdzenie~\ref{thm:distribution-estimate19})
istnieją $C_L>0$ i $m\in\N_0$ takie, że
\[
 |F(f)|\le C_L p_{L,m}(\chi f)
 \le C_L 2^m p_{L,m}(\chi)p_{L,m}(f).
\]
Ostatnia nierówność wynika z reguły Leibniza i tożsamości
$\sum_{\gamma\le\alpha}\binom{\alpha}{\gamma}=2^{|\alpha|}$.
Stąd $F$ jest liniowym funkcjonałem ciągłym.
Wartość nie zależy od wyboru $\chi$: dla dwóch odcięć
$\chi_1,\chi_2$ funkcja $(\chi_1-\chi_2)f$ jest testowa
i znika w otoczeniu $K$. Jej nośnik jest więc zwartym
podzbiorem $\Omega\setminus K$. Z definicji nośnika
i lokalności dystrybucji wynika, że $T$ ją anihiluje. Jeżeli $f\in\mathcal D(\Omega)$, to
$T(f)=T(\chi f)$, ponieważ $(1-\chi)f$ znika w
otoczeniu $K$. Jednoznaczność przedłużenia wynika
z gęstości $\mathcal D(\Omega)$ w $\mathcal E(\Omega)$.
Otrzymana identyfikacja jest zgodna z przedłużeniem
z Lematu~\ref{lem:compact-action21} oraz Uwagą~
\ref{rem:Eprime21}; nie nadajemy obrazowi automatycznie
topologii podprzestrzeni $\mathcal D'$.
\end{proof}
